% Module 3: Cognitive Reasoning -- course essay
% Compile with XeLaTeX: latexmk -xelatex report.tex
\documentclass{article}

\PassOptionsToPackage{numbers,compress}{natbib}
\usepackage[final]{template22}
% Common packages
\usepackage[utf8]{inputenc} % allow utf-8 input
\usepackage[T1]{fontenc}    % use 8-bit T1 fonts
\usepackage{microtype}
\usepackage{times}
\usepackage{graphicx}
\usepackage{amsmath,amssymb,mathbbol}
% \usepackage{algorithmic}
% \usepackage[linesnumbered,ruled,vlined]{algorithm2e}
\usepackage{acronym}
\usepackage{enumitem}
\usepackage[pagebackref=true,breaklinks=true,colorlinks]{hyperref}
\usepackage{balance}
\usepackage{xspace}
\usepackage{setspace}
\usepackage[skip=3pt,font=small]{subcaption}
\usepackage[skip=3pt,font=small]{caption}
\usepackage[dvipsnames]{xcolor}
\usepackage[capitalise]{cleveref}
\usepackage{booktabs,tabularx,colortbl,multirow,array,makecell}
% \usepackage{overpic,wrapfig}

\usepackage{fancyhdr}
\hypersetup{pdfencoding=auto,colorlinks=true,allcolors=black}
\renewcommand{\headrulewidth}{0.5pt}
\renewcommand{\footrulewidth}{0pt}
\fancyhf{}
\fancyhead[C]{}
\fancyhead[C]{}
\fancyfoot[C]{\thepage}

% Handy shorthand
\makeatletter
\DeclareRobustCommand\onedot{\futurelet\@let@token\@onedot}
\def\@onedot{\ifx\@let@token.\else.\null\fi\xspace}
\def\eg{\emph{e.g}\onedot} 
\def\Eg{\emph{E.g}\onedot}
\def\ie{\emph{i.e}\onedot} 
\def\Ie{\emph{I.e}\onedot}
\def\cf{\emph{c.f}\onedot} 
\def\Cf{\emph{C.f}\onedot}
\def\etc{\emph{etc}\onedot} 
\def\vs{\emph{vs}\onedot}
\def\wrt{w.r.t\onedot} 
\def\dof{d.o.f\onedot}
\def\etal{\emph{et al}\onedot}
\makeatother

\definecolor{gray}{gray}{0.9}

% Handy math ops
\DeclareMathOperator*{\argmax}{arg\,max}
\DeclareMathOperator*{\argmin}{arg\,min}
\newcommand{\norm}[1]{\left\Vert #1 \right\Vert}

% % Spacing
\frenchspacing
% \medmuskip=2mu   % reduce spacing around binary operators
% \thickmuskip=3mu % reduce spacing around relational operators
% \setlength{\abovedisplayskip}{3pt}
% \setlength{\belowdisplayskip}{3pt}
% \setlength{\abovecaptionskip}{3pt}
% \setlength{\belowcaptionskip}{3pt}
\setlength\floatsep{0.5\baselineskip plus 3pt minus 2pt}
\setlength\textfloatsep{0.5\baselineskip plus 3pt minus 2pt}
\setlength\dbltextfloatsep{0.5\baselineskip plus 3pt minus 2pt}
\setlength\intextsep{0.5\baselineskip plus 3pt minus 2pt}

\makeatletter
\renewcommand{\paragraph}{%
  \@startsection{paragraph}{4}%
  {\z@}{0ex \@plus 0ex \@minus 0ex}{-1em}%
  {\hskip\parindent\normalfont\normalsize\bfseries}%
}
\makeatother

% Graphics path
\graphicspath{{figures/}}

% Clever references
\crefname{algorithm}{Alg.}{Algs.}
\Crefname{algorithm}{Algorithm}{Algorithms}
\crefname{section}{Sec.}{Secs.}
\Crefname{section}{Section}{Sections}
\crefname{table}{Tab.}{Tabs.}
\Crefname{table}{Table}{Tables}
\crefname{figure}{Fig.}{Fig.}
\Crefname{figure}{Figure}{Figure}

% Acronym
\acrodef{pku}[PKU]{Peking University}

\title{Learning from Surprise:\\
Prediction Error as a Functional Bridge Between Adaptive Networks and Infant Exploration}

\author{%
  Yuanze Lei \quad 2510124204\\
  Peking University
}

\begin{document}
\maketitle

\begin{abstract}
Prediction error---the gap between what we expect and what actually happens---plays a
central role in two very different literatures. In the theory of adaptive networks, Sutton
and Barto (1981) used it to explain how a simple learning element can anticipate future
reward. In developmental psychology, Stahl and Feigenson (2015) showed that 11-month-old
infants learn more about an object when it violates their expectations. This essay compares
the two papers, critically examines the infant experiments, and asks what AI should take
from them. I argue that AI should borrow the functional principles behind biological
learning---prediction-error-driven learning, selective curiosity, and structured
priors---without trying to copy the brain's hardware.
\end{abstract}

\section{Introduction}\label{sec:intro}

Modern AI is very good at recognising patterns, generating text, and playing games. It is
much worse at the questions a toddler answers every day: which parts of experience matter,
when to change one's mind, and how to use old knowledge in a new situation. Large models
learn from far more data than a child ever sees, yet they still struggle with continual
learning, generalisation, and efficient exploration.

Two research traditions have looked at these questions from opposite sides. Adaptive-network
theory asks how simple learning units can pick up predictive structure from experience.
Developmental psychology asks how infants choose what to learn from an overwhelming world.
The two fields use different methods and words, but they have reached a similar conclusion:
\emph{violations of expectation are valuable learning signals}. The conclusion is
functional: it is about the role of prediction error in a learning system, not about the
physical material that implements that role.

This essay has four parts. Section~\ref{sec:history} traces Sutton and Barto's (1981) idea
from animal conditioning to modern reinforcement learning. Section~\ref{sec:critique}
examines Stahl and Feigenson's (2015) four experiments and asks how strong their evidence
really is. Section~\ref{sec:bio} argues that AI should have a \emph{functional} biological
basis and derives three design principles. Section~\ref{sec:reflection} gives my own
reflections.

\section{Historical Context and Modern Applications}\label{sec:history}

\subsection{From ``happening together'' to ``predicting what comes next''}

An early view of learning was based on contiguity: when two events happen together, they
become associated. Hebb's famous rule says that neurons which fire together strengthen
their connection \citep{hebb1949organization}. The rule is simple, but it cannot explain
some basic findings. Suppose one cue already predicts food perfectly. A second cue that
appears at the same time may learn very little---an effect called blocking. Contiguity
alone does not explain why some cues matter and others do not.

The Rescorla--Wagner model gave an answer \citep{rescorla1972theory}. Learning depends on
the gap between the outcome that actually occurs and the outcome predicted by all
available cues:
\begin{equation}\label{eq:rw}
  \Delta V_i = \alpha_i \beta \left(\lambda - \sum_j V_j\right).
\end{equation}
Here $\lambda$ is the size of the outcome, and $\sum_j V_j$ is the total prediction. If the
outcome is fully predicted, the gap is zero and nothing is learned. If the outcome is
better than expected, learning is positive; if it is worse, learning is negative
(extinction). The prediction error is the teacher.

Sutton and Barto's \emph{Toward a Modern Theory of Adaptive Networks: Expectation and
Prediction} (1981) added time to this picture \citep{sutton1981toward}. Their ``adaptive
element'' is a single unit that mimics conditioning, but it is built to match animal
learning more closely than earlier network units. The key feature is that conditioning is
\emph{predictive}: an animal does not simply react to a stimulus; it anticipates what comes
next. The element learns to increase its response rate before the reward arrives.

Three points stand out. First, what is learned is a prediction, not a reaction to the
present. Second, the learning signal is a gap between prediction and outcome, measured at
several points in time. Third, computer simulations showed that the element becomes
sensitive to the most reliable, non-redundant, and earliest predictors of reward. The model
also avoided stability and saturation problems that had troubled earlier networks. The
paper ends by linking the model to synaptic physiology, placing it at the border of
behaviour, computation, and biology.

\subsection{From one element to an industry}

The functional core of the 1981 model---learning from bootstrapped prediction
error---became central to reinforcement learning. In temporal-difference (TD) learning,
the value of a state is updated using the difference between the current prediction and a
better-informed prediction made one step later:
\begin{equation}\label{eq:td}
  \delta_t = r_{t+1} + \gamma V(s_{t+1}) - V(s_t), \qquad
  V(s_t) \leftarrow V(s_t) + \alpha \delta_t .
\end{equation}
In plain words: the next moment gives a slightly better guess than the current moment, and
the difference is used to correct the current guess. This is the ``prediction of
prediction'' idea from 1981, later formalised as TD($\lambda$)
\citep{sutton1988learning} and combined with deep networks in DQN \citep{mnih2015human}.

A striking biological result supports this picture. Schultz, Dayan, and Montague found
that dopamine neurons in the midbrain show the signature of a reward prediction error:
they fire more when reward is better than expected, not at all when reward is fully
predicted, and less when an expected reward is omitted \citep{schultz1997neural}. The
biology is far more complex than a single number, but the computational role is the same.

Modern systems keep the function but change the implementation. Three differences matter.
First, the \emph{target} has changed: Sutton and Barto tried to explain animal
conditioning, while modern reinforcement learning optimises an engineered objective. In
language models trained with human feedback, the reward comes from a learned preference
model \citep{ouyang2022training}, not from food or pain. Second, the \emph{scale} has
changed: a single element with scalar predictions has become deep networks with replay
buffers, target networks, and many parallel actors \citep{mnih2015human}. Third,
\emph{biological plausibility} has stopped being a design constraint: the 1981 paper links
its element to synaptic biochemistry, whereas modern systems use global backpropagation
and centralised training, neither of which is a realistic model of the brain.

What survives is not a piece of anatomy but a principle: a learner should use the gap
between prediction and outcome to update its model, and it should propagate that gap
backwards in time. This is the first piece of evidence for the thesis in
Section~\ref{sec:bio}: biological inspiration can be useful even when the biological
material is left behind.

\section{Critical Analysis of Modern Research}\label{sec:critique}

\subsection{What the four experiments actually did}

Stahl and Feigenson's \emph{Observing the Unexpected Enhances Infants' Learning and
Exploration} (2015) starts from a simple problem: a baby cannot learn everything, so how
does she decide what to learn? Their answer is that violations of expectation act as a
filter. Infants already have ``core knowledge'' about objects---for example, that solid
objects cannot pass through walls, and that objects need support
\citep{spelke2007core}. When an event breaks one of these rules, the baby may treat it as a
special chance to learn.

Because infants cannot speak, the researchers used two measures: where the baby looked,
and how the baby played. Across four experiments ($N=110$), all participants were 11
months old.

Experiment~1 ($N=40$) had four groups. In the solidity event, a toy car rolled down a ramp
toward a wall. Half the infants saw the car stop at the wall (consistent); the other half
saw it pass through (violation). In the continuity event, a ball was hidden behind the left
screen; when both screens were lifted, the ball was either still behind the left screen
(consistent) or behind the right screen (violation). Every infant saw the outcome for only
10 seconds.

Next, the experimenter taught the infant a new property of the object: when the object was
moved, a squeaking sound played from a hidden speaker. Finally, the infant saw the target
object and a new distractor object. In the baseline phase, both objects were still. In the
test phase, both objects moved and the taught sound played. If the infant had learned
``this object makes this sound'', she should look more at the target when the sound played.
The learning score was the proportion of looking at the target during the test minus the
proportion during the baseline.

The violation groups learned (solidity: $t(9)=3.09$, $p=.01$; continuity: $t(9)=3.72$,
$p=.005$). The consistent groups did not ($p=.31$ and $p=.14$). The overall difference was
significant ($F(1,36)=10.69$, $p=.002$, partial $\eta^2=.23$).

Experiment~2 ($N=20$) repeated the violation condition but played a new sound at test. The
learning score was close to zero ($p=.18$ and $p=.97$). The authors concluded that the
effect in Experiment~1 was about learning a specific object--sound link, not a general
preference for the surprising object.

Experiment~3 ($N=10$) repeated the continuity violation but taught a new object that had
not been part of the violation. The infants did not learn ($p=.94$), and this differed from
Experiment~1 ($p=.048$). The authors concluded that the effect was specific to the object
that had broken expectations, not a general increase in arousal.

Experiment~4 ($N=40$) looked at exploration. Infants saw either a solidity violation (car
through wall) or a support violation (car floating beyond the table edge). They then
played freely with the target object and a new distractor for 60 seconds. Infants in the
violation conditions explored the target more ($p=.02$). Their actions also matched the
type of violation: after a solidity violation they banged the object more ($p=.029$);
after a support violation they dropped it more ($p=.056$, marginal). The same pattern did
not appear for the distractor object ($p=.80$).

\subsection{How strong is the evidence?}

The study has real strengths. The events were matched in every respect except the outcome,
and every infant saw the outcome for exactly 10 seconds. There was no difference between
conditions in looking time at the outcome ($F(1,36)=0.002$, $p=.96$), so the result cannot
be explained by longer exposure. The hidden property was genuinely new, so the test
measured learning rather than recognition. Experiments~2 and~3 were designed as controls,
and the double dissociation in Experiment~4 is a strong sign of specificity.

But the evidence is not as clean as the paper's language suggests. The weakest point is
Experiment~2. The authors argue that if infants simply preferred the surprising object,
they should look at it in Experiment~2 as well, because the object was the same. But the
learning score is a difference score: test minus baseline. If an infant likes the object
equally in both phases, the score is zero whether or not any learning occurred. A zero
score therefore cannot separate ``no learning'' from ``a constant preference for the
object''. Experiment~2 provides some evidence against a general visual preference, but it
does not settle the question.

Experiment~1 does more work here, because it subtracts the baseline. The violation groups
showed a larger increase from baseline to test than the consistent groups. Still, the
baseline (5 seconds, objects still) and the test (10 seconds, objects moving, sound
playing) are not perfectly comparable. A ``surprise + movement'' effect could produce part
of the increase.

Other limits are worth noting. The samples are small (ten infants per cell), and a
non-significant result in Experiment~2 is not proof of no effect. The ``learning'' is an
arbitrary object--sound mapping, not an understanding of solidity or support. All
participants were 11 months old and tested in a laboratory, so the findings may not
generalise across ages, cultures, or species. Effect sizes are moderate, and one key
contrast in Experiment~4 is only marginal ($p=.056$).

Overall, the hypothesis--design--conclusion chain is largely consistent. The study
convincingly shows that violations change what infants learn and how they explore. It does
not fully separate learning from attention, and its control for visual preference is
weaker than the paper claims.

\subsection{Why this matters for AI}

The infant result does not mean ``maximise surprise''. It means that a violation should
identify \emph{which object} and \emph{which hypothesis} deserve attention. This matters
for AI, where raw surprise is easy to exploit: an agent rewarded for prediction error may
stare at a noisy television instead of learning about the world. Curiosity-driven methods
such as self-supervised prediction \citep{pathak2017curiosity} try to avoid this, but the
deeper lesson is that curiosity should be gated by relevance and learnability, not by the
size of the error alone.

There is also a warning for the analogy. Infants' expectations come from core knowledge
that is present before the experiment. An AI's expectations come from its training data.
A ``violation'' for an AI may be a distribution shift, not a physical anomaly. The
comparison is useful only if we treat it as a hypothesis about function, not as an identity
between babies and machines.

\section{Biological Foundations in AI Development}\label{sec:bio}

\subsection{The thesis}

My thesis is that future AI systems \emph{should} have a biological basis, but the right
basis is functional rather than anatomical. AI should borrow the computational roles that
biological systems give to prediction error, surprise, and structured priors. It does not
need to copy neurons, synapses, or brain regions.

The difference matters because ``biological basis'' can mean two things. At the
implementation level, it means spiking neurons, synaptic plasticity, and detailed anatomy.
At the functional level, it means principles such as error-driven learning, selective
attention, curiosity, compositionality, and hierarchical control. The first is a research
programme in its own right, but it is not a requirement for better learning systems. The
second is more tractable and more directly supported by the two papers.

\subsection{Three design principles}

\paragraph{P1: Prediction error should be a first-class learning signal.}
The Rescorla--Wagner model, the 1981 element, and TD learning all replace a passive
association with an active prediction, and use the error of that prediction as the teacher.
Dopamine data support this at the biological level \citep{schultz1997neural}, and
self-supervised prediction supports it in AI \citep{pathak2017curiosity}. The design
implication is simple: train systems not only on external rewards or labels, but also on
auxiliary predictions about what will happen next. World models, model-based planning, and
self-supervised prediction are natural examples.

\paragraph{P2: Surprise should select what to learn, and learnability should gate the choice.}
The infants in Stahl and Feigenson's study do not explore everything after a violation.
They explore the object that broke their expectations, and they perform actions that match
the type of violation. This is a scheduling policy, not a single curiosity score.
Intrinsic motivation has a long history in AI \citep{schmidhuber2010formal}, but raw
prediction error can be gamed. A better translation uses surprise to propose hypotheses,
then allocates effort according to expected information gain, learning progress, or
controllability. The key word is \emph{targeted}.

\paragraph{P3: Structured, compositional priors should constrain the hypothesis space.}
Core knowledge gives infants a small set of abstract expectations about objects, agents,
and number \citep{spelke2007core}. These expectations do not tell infants what to believe;
they tell them which hypotheses are worth considering. The AI analogue is not a hand-coded
knowledge base. It is a set of architectural biases: object-centric representations,
modular and hierarchical networks, and compositional structures that support systematic
generalisation. Lake and colleagues make a related argument when they propose that
human-like machines should build causal world models, ground learning in intuitive
theories, and exploit compositionality and learning-to-learn \citep{lake2017building}.

\subsection{Why not copy the brain?}

Three reasons favour functional grounding. First, we still do not understand biological
learning in enough detail; building around a simplified synaptic story risks turning a
metaphor into a design constraint. Second, engineering matters: global backpropagation,
large-scale parallelism, and centralised training work extremely well even though they are
not biologically plausible. Third, the history of TD learning is a proof of concept. It
began as a model of animal conditioning, was supported by dopamine recordings, and then
became central to systems that make no claim to biological realism. The principle survived
because it could be separated from its substrate.

The obvious objection is that scale alone may produce these abilities. This is a serious
position, and it should not be dismissed. But scale buys competence; it does not
automatically buy sample-efficient exploration, continual learning, or causal
understanding. The question is empirical: do biologically inspired inductive biases
reduce the sample complexity of the abilities we care about? The debate should be settled
with benchmarks that isolate each principle---for example, comparing learnability-gated
curiosity with raw surprise on noisy-TV and continual-learning tasks.

\section{Personal Insights and Reflections}\label{sec:reflection}

What I find most striking is that both papers treat surprise as a \emph{selection}
mechanism rather than an error to be suppressed. In most machine learning, surprise is a
symptom of failure. In both literatures, it is also an opportunity: it marks the part of
the world where the current model is wrong and where learning is likely to be most
valuable. That reframing turns a scalar loss into a policy over what to investigate next.

This matters for language models. Pre-training is token-level prediction-error
minimisation, and reinforcement learning from human feedback adds a preference-based
reward \citep{ouyang2022training}. But current systems have no explicit mechanism for using
surprise to decide when to ask a question, run an experiment, or update an internal model.
A model that could notice when its prediction is unreliable, generate a targeted probe,
and incorporate the result would be closer to the infant in Stahl and Feigenson's study
than a model that simply scales up next-token prediction.

I am also struck by the risks. It is easy to anthropomorphise the infant results and assume
that an agent with a curiosity reward is doing what a baby does. The analogy hides
important differences: infants have core knowledge, bodies, and a developmental history;
current agents have training distributions, reward models, and offline datasets. If
surprise is optimised directly, agents may seek unpredictable noise rather than learnable
structure. The infant evidence suggests a corrective: the value of a violation lies not in
its novelty, but in its relevance to a specific hypothesis.

My own view is cautiously optimistic about functional biological grounding and sceptical
about anatomical imitation. Cognitive science is most useful to AI when it formulates the
right questions---what to learn, when to revise, how to generalise---rather than when it
supplies a blueprint for implementation.

\section{Conclusion}\label{sec:conclusion}

Prediction error connects an adaptive element designed to explain animal conditioning, a
temporal-difference algorithm that became a cornerstone of reinforcement learning, and an
infant who learns more about an object that has just broken her expectations. The two
papers are separated by 34 years, two disciplines, and very different methods, yet they
reach the same functional claim: a learner should treat violations of expectation as
invitations to learn.

AI should take that claim seriously, but not literally. It does not need spiking neurons or
a copy of the brain. It needs learning systems whose errors drive representation learning,
whose surprise is selective and learnability-gated, and whose representations are
structured enough to support composition and generalisation. Biology supplies the
principles; engineering supplies the rest.

\bibliographystyle{plainnat}
\bibliography{reference}

\end{document}
